\chapter{Análise dos Resultados}
\label{cap:analise}

Este capítulo apresenta a interpretação crítica e a discussão aprofundada dos achados empíricos e técnicos obtidos ao longo da pesquisa, respondendo ao quinto objetivo específico do estudo (OE5). Inicialmente, a Seção \ref{sec:triangulacao_resultados} interpreta os indicadores estatísticos coletados em campo e realiza a triangulação metodológica entre as frentes quantitativa, qualitativa e de telemetria (atividade A5.2). Na sequência, a Seção \ref{sec:analise_comparativa} desenvolve a análise comparativa multicritério entre a prova de conceito original e a nova versão do ElicitBot (atividade A5.3). Por fim, a Seção \ref{sec:contribuicoes} sintetiza as contribuições acadêmicas e práticas da pesquisa frente à literatura, enquanto a Seção \ref{sec:limitacoes_estudo} delimita as limitações e ameaças à validade do trabalho.

\section{Interpretação dos Dados e Triangulação Metodológica}
\label{sec:triangulacao_resultados}

A avaliação empírica do ElicitBot junto aos $N = 65$ estudantes de Engenharia de Software foi conduzida sob uma abordagem mista (quanti-qualitativa), confrontando os indicadores métricos de desempenho do agente e os registros de telemetria do servidor com a experiência concreta de uso reportada nas questões abertas do instrumento.

Sob a perspectiva quantitativa, o estudo revelou uma consistência de topo expressiva: todos os 20 critérios avaliados apresentaram tanto Mediana ($Md = 5,0$) quanto Moda ($Mo = 5,0$) iguais à nota máxima da escala Likert, alcançando média global de $\mu = 4,49$. Em 9 dos 20 itens (45\% do instrumento), incluindo \textit{Adesão ao Fluxo Lógico}, \textit{Investigação do Problema Real}, \textit{Adesão às Restrições}, \textit{Fidelidade de Regras}, \textit{Verificabilidade} e \textit{Rastreabilidade de Requisitos}, não houve qualquer registro de avaliação discordante (notas 1 ou 2 iguais a zero). Ademais, o equilíbrio nas avaliações percebidas entre as dimensões de modelagem do negócio (D2, com $\mu = 4,42$), concepção de soluções (D3, com $\mu = 4,50$) e especificação de requisitos (D4, com $\mu = 4,50$) indica consistência na percepção dos estudantes ao longo das diferentes etapas do fluxo de elicitação.

Para interpretar adequadamente o cruzamento entre esses indicadores, a telemetria e os relatos discentes, torna-se fundamental considerar uma condição estrutural da dinâmica aplicada em sala de aula: as equipes já possuíam o escopo e o levantamento de requisitos de seus projetos concebidos previamente. Dessa forma, em vez de interagirem sob o ritmo orgânico e gradual de uma elicitação exploratória (na qual o usuário descobre e amadurece o problema junto ao assistente ao longo de várias sessões), os estudantes utilizaram a ferramenta para transpor e validar um projeto já existente dentro de uma única sessão de 90 minutos. A triangulação integrada demonstra como esse comportamento de transposição acelerada explica os fenômenos observados nos cinco eixos analíticos a seguir.

\subsection{Eficácia da Máquina de Estados frente à Rigidez Metodológica}

O resultado quantitativo mais expressivo da avaliação foi a liderança da dimensão de Condução Conversacional e Controle de Processo (D1, com $\mu = 4,58$), impulsionada por 100,0\% de aprovação positiva (\textit{Top-2 Box}) em \textit{Adesão ao Fluxo Lógico} ($\mu = 4,78$ e $\sigma = 0,41$), 96,9\% em \textit{Validação de Transição} ($\mu = 4,83$) e 93,8\% em \textit{Investigação do Problema Real} ($\mu = 4,68$). Esses indicadores de percepção sugerem que a Máquina de Estados implementada no \textit{System Prompt} foi reconhecida pelos estudantes como um mecanismo consistente para conduzir a entrevista de forma progressiva e conter o salto prematuro para funcionalidades, mesmo diante de grupos que já detinham a lista final de requisitos em mãos.

Na análise qualitativa, esse mesmo rigor estrutural manifestou-se nos relatos classificados como Limitações Aceitáveis e Delimitação de Escopo. Os estudantes apontaram a recusa do agente em migrar para abordagens alternativas a Histórias de Usuário, como o \textit{Feature-Driven Development} (\texttt{[P39, P40, P47]}), a manutenção do aluno no papel de \textit{stakeholder} durante a entrevista consultiva (\texttt{[P5, P8]}) e a capacidade do assistente de responder brevemente a provocações fora de domínio (como receitas culinárias) e reconduzir imediatamente o diálogo ao escopo técnico (\texttt{[P39, P43, P47, P52]}). Essa convergência entre notas e relatos indica que as travas metodológicas foram percebidas pelos avaliadores como um mecanismo consistente de disciplina analítica e aderência estrutural, refletida nos 90,2\% de aprovação percebida em \textit{Adesão ao Formato Padrão} ($\mu = 4,57$).

\subsection{Atravessamento Acelerado do Funil e Gargalos nas Etapas Finais}

A análise do funil de progressão autodeclarado no questionário demonstrou que 92,3\% dos participantes (60 de 65 estudantes) percorreram todas as etapas de entendimento do domínio, modelagem do negócio e concepção da solução, alcançando a fase de \textit{Requisitos Funcionais} (Fase 10) em uma única aula. Enquanto 41,5\% da amostra (27 estudantes) relatou ter atingido a última etapa de \textit{Confirmar Escopo} (Fase 12), outros 50,8\% (33 estudantes) encerraram o tempo de laboratório concentrados nas Fases 10 (38,5\%) e 11 (12,3\%).

Esse avanço acelerado pelas fases iniciais e a posterior concentração nas Fases 10 e 11 foram confirmados pela telemetria da entidade \texttt{MetricaLLM}. Como as equipes já conheciam previamente o contexto, os problemas e os atores de seus projetos, as seis etapas iniciais de domínio foram percorridas com rapidez (frequentemente mediante respostas curtas de aprovação, o que explica a mediana de apenas 19,0 \textit{tokens} nas mensagens dos usuários). Consequentemente, 211 dos 519 turnos totais da sessão (40,7\%) ocorreram com o atributo \texttt{Fase do Cache} já estabilizado em \texttt{validacao\_solucao} (equipes atuando nas Fases 10 e 11), e outros 21 turnos (4,0\%) ocorreram com o \textit{cache} em \texttt{validacao\_requisitos} (equipes atuando na Fase 12). O cruzamento com a triagem qualitativa explica por que o ritmo de avanço desacelerou justamente nas Fases 10 e 11:

\begin{itemize}
    \item \textbf{Fator Metodológico de Transposição (Limitação Aceitável \texttt{LA-06}):} Ao chegarem à Fase 10 com uma lista prévia de requisitos já elaborada para a disciplina, parte das equipes tentou despejar múltiplos requisitos simultaneamente (registrando-se um turno em que o modelo gerou 10.024 \textit{tokens} JSON de saída de uma só vez). A diretriz de interação incremental e validação unitária para requisitos com pendências (relatada por 3 participantes em \texttt{LA-06}) gerou fricção com as tentativas de inserção/validação massiva em lotes (\textit{batches}), exigindo que cada grupo revisasse regras de negócio, critérios de aceitação e vínculos de rastreabilidade pendentes requisito a requisito.
    \item \textbf{Fator Sistêmico (Correção Imediata \texttt{CI-02}):} A pressa em corrigir e aprovar requisitos acumulados na Fase 11 acentuou a ocorrência da falha de concorrência (\textit{Validação de Requisitos}, relatada por 14 participantes), na qual a tentativa do LLM de atualizar pendências de um requisito e submetê-lo à validação no mesmo turno ocasionou erros de submissão e interrupções pontuais, atrasando a chegada de parte dos grupos à Fase 12.
\end{itemize}

Mesmo diante dessas fricções nas etapas finais, os estudantes mantiveram uma percepção positiva quanto aos artefatos gerados na Dimensão D4 ($\mu = 4,50$), com destaque para 96,6\% de aprovação percebida em \textit{Rastreabilidade de Requisitos} ($\mu = 4,60$).

\subsection{Persistência Estruturada e \textit{Human-in-the-Loop} na Retenção de Contexto}

A leitura integrada também elucida o comportamento do único item que apresentou média inferior a 4,30 no instrumento: o critério \textit{Retenção de Contexto} (pertencente à Dimensão D1), que registrou média $\mu = 4,09$, desvio padrão $\sigma = 1,17$ e taxa de aprovação positiva de 69,2\% (atuando como um valor atípico isolado que reduziu a média de D1 de $\mu = 4,67$ para $\mu = 4,58$, embora preservando $Md = 5,0$ com 53,8\% de notas máximas). A comparação com os índices das fases subsequentes revela que essa percepção pontual não comprometeu a integridade semântica dos projetos, visto que \textit{Adesão às Restrições} atingiu 96,8\% de aprovação ($\mu = 4,63$) e \textit{Alinhamento da Solução} alcançou 95,3\% ($\mu = 4,53$).

A triangulação entre a telemetria e os relatos abertos demonstra que a percepção de perda de contexto decorreu diretamente do choque entre o envio massivo de documentação prévia pelos alunos e a arquitetura de memória do sistema:

\begin{enumerate}
    \item \textbf{Despejo Prematuro de Documentação (\textit{Information Dumping}) vs. Janela Deslizante:} A telemetria revelou que algumas equipes colaram documentos inteiros do projeto em um único turno (com mensagens de usuário atingindo 3.691 \textit{tokens} logo no primeiro turno e até 39.054 \textit{tokens} em etapas intermediárias). Quando o estudante despeja especificações completas nas fases iniciais, a trava de \textit{Defesa do Escopo} impede a criação precoce de requisitos e retém apenas apontamentos sintéticos em \texttt{topicos\_em\_espera}. Como o histórico bruto opera sob uma janela deslizante de 20 mensagens (10 turnos) combinada ao modelo leve \textit{Gemini 3.1 Flash-Lite}, o texto bruto colado prematuramente deixa de compor a janela ativa após 10 turnos, fazendo com que detalhes não convertidos em artefatos precisem ser reinformados no futuro.
    \item \textbf{Rigor de Validação diante de Entradas Vagas:} Conforme relatado por participantes na fase de \textit{Objetivos de Negócio} (\texttt{[P34, P36]}), quando o usuário fornece respostas superficiais que não satisfazem os critérios semânticos da etapa, o agente reformula a pergunta para aprofundar o levantamento, comportamento que pode ser interpretado por parte dos alunos como esquecimento de informações já enviadas.
    \item \textbf{Autonomia de Refinamento Manual (\textit{Human-in-the-Loop}):} Como os artefatos consolidados são persistidos no PostgreSQL e injetados na Memória Estável a cada turno, a estrutura central do projeto permaneceu intacta, e a barra lateral com suporte a edição manual foi utilizada por 12 participantes (\texttt{LA-01}) para ajustar descrições ou títulos genéricos, restando apenas ajustes pontuais de identidade por \textit{Tag} e sincronização de exclusão, já sanados nas correções imediatas (\texttt{CI-04}).
\end{enumerate}

\subsection{Fricções de Interface, Falhas Operacionais e Percepção de Satisfação e Adoção}

O cruzamento entre o perfil de interação das equipes e os relatos qualitativos evidencia que os incidentes reportados abrangeram tanto aspectos ergonômicos da interface quanto falhas que impactaram tarefas centrais da elicitação. No plano da interface, a concentração de apontamentos na caixa de entrada de texto (\texttt{CI-01}, com 22 relatos) e a sugestão de ingestão de arquivos (\texttt{TF-01}, com 10 relatos) relacionam-se diretamente ao contexto de transposição de projetos pré-documentados: ao colar blocos extensos de texto no \textit{chat}, a limitação de um campo \texttt{<input>} de linha única sem quebra via \texttt{Shift + Enter} tornou-se recorrente, somando-se a 5 relatos sobre rolagem manual e fechamento restrito de modais (\texttt{CI-03}). No plano funcional e de estado, a falha de concorrência na Fase 11 (\texttt{CI-02}, com 14 relatos) atingiu uma atividade central do processo (a validação de requisitos), provocando erros de submissão e interrupções que atrasaram a progressão de parte das equipes, enquanto ocorrências pontuais de duplicação semântica e menção a \textit{Tags} descartadas (\texttt{CI-04}, com 5 relatos) exigiram intervenções corretivas no \textit{prompt}.

Em resposta a esses impactos observados em sala de aula, as quatro frentes de correção imediata (\texttt{CI-01} a \texttt{CI-04}) foram implementadas e verificadas funcionalmente na atividade A4.4, embora sua eficácia em uso coletivo ainda demande nova avaliação com estudantes. Paralelamente, os indicadores da Dimensão D5 ($\mu = 4,39$) mostram que, mesmo com as limitações vivenciadas na sessão avaliada, a percepção global reportada no questionário manteve-se predominantemente favorável: 90,8\% avaliaram positivamente a \textit{Interação e Resultados} ($\mu = 4,40$ e $Md = 5,0$, com um registro discordante) e 83,1\% indicaram concordância quanto ao \textit{Potencial de Uso} do ElicitBot em projetos acadêmicos ou profissionais futuros ($\mu = 4,38$ e $Md = 5,0$, com 58,5\% de notas máximas). Esses achados indicam que o valor percebido da condução estruturada coexistiu com gargalos que precisaram ser endereçados para aprimorar a experiência de uso.

\subsection{Desempenho Operacional e \textit{Cloud Economics} sob Compressão Longitudinal}

Por fim, os registros da entidade \texttt{MetricaLLM} ($N = 519$ turnos) permitiram avaliar o comportamento do modelo financeiro (\textit{Cloud Economics}) sob um cenário ainda mais exigente do que o previsto na Fase 1. No planejamento inicial, projetou-se que um projeto cresceria gradualmente ao longo de quatro sessões semanais de 25 turnos cada. Na prática do UAT, devido ao conhecimento prévio do escopo pelos grupos, as equipes comprimiram o equivalente a todo o ciclo de vida do projeto dentro de uma única sessão de 85 minutos (média de 37,1 turnos por equipe), fazendo com que a Memória Estável (até 6.221 \textit{tokens}) e a Memória Dinâmica (até 17.251 \textit{tokens}) atingissem o topo de volume de dados na mesma aula.

Mesmo sob essa compressão longitudinal e com um \textit{System Prompt} de 9.008 \textit{tokens} (16 vezes superior aos 550 \textit{tokens} estimados preliminarmente na Fase 1), a latência mediana manteve-se em 5.143 ms (P95 de 8.580 ms) e o faturamento total do dia no provedor Google Cloud foi de apenas \textbf{R\$ 20,84} para atender os 78 estudantes presentes. A decomposição entre o custo de processamento apurado na tabela \texttt{MetricaLLM} (US\$ 1,82, ou cerca de R\$ 10,38) e a fatura bruta do provedor (R\$ 20,84) evidencia dois efeitos diretos dessa dinâmica acelerada:

\begin{itemize}
    \item \textbf{Eficiência de Leitura do \textit{Context Caching}:} A Memória Estável garantiu uma taxa de acertos de \textit{cache} (\textit{cache hit}) em 74,0\% das requisições, aplicando o desconto de 90\% sobre 58,7\% de todos os \textit{tokens} de entrada enviados no dia e mantendo o custo médio de processamento em US\$ 0,0035 por turno.
    \item \textbf{Sobreposição de Armazenamento Temporal (\textit{Storage TTL}) por Avanço Rápido de Fases:} Em um ritmo normal de descoberta, uma equipe permaneceria em uma mesma fase por mais tempo, gerando menos objetos de \textit{cache}. Como as equipes já tinham o projeto pronto e avançaram rapidamente pelas 12 etapas, ocorreram 135 eventos de recriação de \textit{cache} (\texttt{Cache Novo = True}, média de 9,6 recriações por equipe, ou um novo objeto a cada 3,8 turnos), somando 1,38 milhão de \textit{tokens} gravados. Como cada objeto criado na Vertex AI retém cobrança de armazenamento temporal (\textit{Storage TTL}, tarifado a US\$ 1,00 por milhão de \textit{tokens}/hora) até expirar seu tempo de vida padrão, múltiplos objetos de fases consecutivas permaneceram simultaneamente ativos na memória da nuvem, adicionando cerca de US\$ 1,38 (aproximadamente R\$ 7,85) à fatura, além de requisições que sofreram reversão transacional (\textit{rollback}) no \textit{backend} durante o erro da Fase 11 (\texttt{CI-02}).
\end{itemize}

Observa-se, portanto, que a sessão experimental configurou um cenário específico de alta densidade de transições em aula única, caracterizado por frequentes invalidações e recriações de \textit{cache} em um intervalo curto. O fato de o custo faturado ter permanecido em R\$ 1,49 por equipe (R\$ 0,27 por estudante presente) sob essas condições oferece um indício favorável de viabilidade financeira no recorte avaliado. Como expectativa a ser verificada em estudos longitudinais futuros, admite-se a hipótese de que, em dinâmicas de descoberta gradual ao longo de várias semanas (com menor taxa horária de recriação de objetos de \textit{cache} e eventuais otimizações de evicção), a parcela associada ao armazenamento temporal (\textit{Storage TTL}) tenda a reduzir-se por sessão, embora o custo acumulado do semestre dependa do volume total de interações e do comportamento de uso de cada turma.

\section{Análise Comparativa entre a Prova de Conceito e a Nova Versão}
\label{sec:analise_comparativa}

Em cumprimento à atividade A5.3, esta seção apresenta a análise comparativa multicritério entre a versão inicial do ElicitBot, desenvolvida como prova de conceito por \citeonline{junior2025}, e a versão evoluída resultante desta pesquisa. Cabe ressaltar que a comparação estabelece uma distinção metodológica entre mudanças estruturais e funcionais diretamente observáveis na implementação das duas versões e resultados empíricos de uso. Como as duas versões atenderam a propósitos distintos e foram avaliadas sob protocolos, instrumentos, domínios de projeto, perfis de interação (sessões individuais frente a equipes com escopo pré-concebido) e participantes diferentes, eventuais diferenças nas percepções reportadas não representam ganhos medidos sob protocolo experimental equivalente e não devem ser atribuídas exclusivamente às modificações implementadas no sistema. Enquanto a prova de conceito teve como finalidade explorar a viabilidade inicial de LLMs em entrevistas simuladas individuais, a nova versão foi projetada para viabilizar a operação concorrente, persistente e controlada da ferramenta em turmas da disciplina de Requisitos de Software.

A Tabela \ref{tab:comparativo_versoes} sintetiza as características estruturais, metodológicas e funcionais de cada versão segundo oito dimensões de análise.

\noindent
{
\footnotesize 
\renewcommand{\arraystretch}{1.15}
\setlength{\LTleft}{0pt}
\setlength{\LTright}{0pt}

\begin{xltabular}{\textwidth}{|>{\raggedright\arraybackslash}p{2.8cm}|>{\raggedright\arraybackslash}X|>{\raggedright\arraybackslash}X|}
    \caption{Comparativo Sistemático entre a Prova de Conceito e a Versão Evoluída do ElicitBot}
    \label{tab:comparativo_versoes} \\
    \toprule
    \textbf{Dimensão Analisada} & \textbf{Prova de Conceito \cite{junior2025}} & \textbf{Versão Evoluída (Este Trabalho)} \\
    \midrule
    \endfirsthead

    \multicolumn{3}{c}%
    {{\bfseries \tablename\ \thetable{} -- continuação da página anterior}} \\
    \toprule
    \textbf{Dimensão Analisada} & \textbf{Prova de Conceito \cite{junior2025}} & \textbf{Versão Evoluída (Este Trabalho)} \\
    \midrule
    \endhead

    \midrule \multicolumn{3}{r}{{Continua na próxima página...}} \\
    \endfoot

    \bottomrule
    \multicolumn{3}{c}{\small Fonte: Elaborada pelo autor.} \\
    \endlastfoot
    
    \textbf{Propósito e Escopo de Operação} & Protótipo exploratório voltado à validação conceitual da elicitação assistida por IA em sessões individuais agendadas. & Plataforma educacional voltada à operação concorrente em sala de aula para turmas de 40 a 50 estudantes simultâneos. \\
    \midrule
    \textbf{Modelo Metodológico de Elicitação} & Diálogo livre orientado diretamente ao levantamento de Requisitos Funcionais, guiado por 10 diretrizes gerais de condução e resumo textual por tópico. & Condução estruturada por uma Máquina de Estados Finita de 12 etapas sequenciais, dividida entre entendimento do domínio (Fases 1 a 6) e definição da solução (Fases 7 a 12). \\
    \midrule
    \textbf{Artefatos de Domínio e Rastreabilidade} & Extração centrada em uma tipologia de artefato (Requisito Funcional composto por descrição, História de Usuário no padrão Connextra, regras de negócio e critérios de aceitação). & Modelagem relacional de oito categorias de artefatos (\texttt{CTX}, \texttt{PROB}, \texttt{ENV}, \texttt{NEC}, \texttt{OBJ}, \texttt{RES}, \texttt{SOL} e \texttt{RF}) interligadas por matriz de rastreabilidade via \textit{Tags} imutáveis e auditadas por dupla validação (servidor e IA). \\
    \midrule
    \textbf{Arquitetura de Software e Hospedagem} & Aplicação monolítica em Python acoplada ao \textit{framework} Streamlit, com execução em \textit{thread} auxiliar local e hospedagem no \textit{Streamlit Community Cloud}. & Arquitetura cliente-servidor desacoplada (interface em React/TypeScript e API RESTful assíncrona em FastAPI/Uvicorn), conteinerizada via Docker com Nginx e hospedada em servidor institucional (CEDIS/UnB). \\
    \midrule
    \textbf{Persistência de Estado e Sessão} & Armazenamento temporário em memória RAM vinculado à sessão ativa do navegador (\texttt{st.\allowbreak session\_\allowbreak state}), com reinicialização manual via botão ``Limpar Conversa''. & Persistência relacional transacional (ACID) em PostgreSQL 16 via SQLAlchemy 2.0 assíncrono e padrão \textit{Read-Release-Write}, suportando múltiplos projetos por usuário e retomada contínua de histórico. \\
    \midrule
    \textbf{Comunicação com LLM e Engenharia de Contexto} & Duas chamadas HTTP por turno (diálogo com reenvio cumulativo do histórico e extração em segundo plano das últimas 4 mensagens via separação de texto livre). & Chamada única por turno com Saídas Estruturadas (\textit{JSON Schema}), combinada ao Gerenciamento Bipartido de Memória (\textit{Context Caching} na Memória Estável, Memória Dinâmica e Janela Deslizante de 20 mensagens). \\
    \midrule
    \textbf{Interface, Controle do Usuário e Exportação} & Barra lateral com guia de instruções e cartões expansíveis de requisitos, permitindo exclusão individual de itens e exportação em planilha Excel (\texttt{.xlsx}). & Painel de gestão de projetos (\textit{Dashboard}), área de trabalho com oito abas categorizadas, edição manual em modal (\textit{Human-in-the-Loop}), reversão transacional de turno, avaliação de utilidade e exportação de DRS em PDF, Markdown e JSON. \\
    \midrule
    \textbf{Governança, Segurança e Observabilidade} & Acesso direto via endereço público e proteção de credenciais da API por variáveis de ambiente (\texttt{.env} e \textit{Streamlit Secrets}). & Autenticação institucional via Microsoft Entra ID (SSO/JWT), controle por papéis (RBAC), gestão de \textit{Whitelist} e janelas letivas, \textit{throttling}, versionamento de \textit{prompts} (\textit{LLMOps}), exportação de telemetria em CSV e console de \textit{logs} em tempo real. \\
\end{xltabular}
}

A análise detalhada das dimensões apresentadas na Tabela \ref{tab:comparativo_versoes} permite caracterizar a evolução técnica e funcional ocorrida entre os dois ciclos de pesquisa:

\begin{enumerate}
    \item \textbf{Transição do Fluxo Livre para a Máquina de Estados Orientada ao Domínio:} Na prova de conceito, o módulo \texttt{LLMManager} conduzia a entrevista a partir da pergunta inicial sobre o que o sistema deveria fazer, extraindo diretamente Requisitos Funcionais a partir do quarto turno de conversa. Na versão evoluída, o processo foi expandido para uma Máquina de Estados de 12 fases que percorre a compreensão do contexto organizacional, problemas, partes interessadas, necessidades, objetivos de negócio e restrições antes de conceber alternativas de solução e derivar requisitos funcionais. Essa ampliação metodológica permitiu que os requisitos gerados na Fase 10 fossem estruturalmente vinculados às dores e necessidades levantadas nas etapas iniciais, além de introduzir uma etapa dedicada de verificação determinística e semântica na Fase 11.
    \item \textbf{Desacoplamento Arquitetural e Persistência Relacional:} A versão original adotou o \textit{framework} Streamlit e o armazenamento em memória de sessão (\texttt{st.\allowbreak session\_\allowbreak state}), escolhas adequadas para reduzir o tempo de prototipação e viabilizar testes exploratórios individuais. Para atender ao requisito de operação simultânea por dezenas de estudantes sem perda de dados por instabilidades na rede sem fio universitária, a nova versão separou a camada de apresentação (executada no navegador do cliente via React) da camada de serviço (FastAPI assíncrono) e introduziu o banco de dados PostgreSQL 16. O gerenciamento de estado passou a contar com um painel de projetos (\textit{Dashboard}), permitindo ao estudante pausar, retomar, renomear ou excluir sessões de elicitação de forma independente.
    \item \textbf{Unificação de Requisições e Orquestração de Memória em Nuvem:} No sistema precursor, a separação entre a geração da resposta conversacional e a rotina de extração (\texttt{\_extract\_requirements}) exigia duas requisições HTTP por interação e realizava a leitura dos campos por meio de fatiamento textual de blocos separados por quebras de linha. Na nova arquitetura, o uso de Saídas Estruturadas (\textit{Structured Outputs}) unificou o diálogo, a transição de estado e os cinco comandos do Motor de Operações em uma única chamada validada por esquemas Pydantic. Adicionalmente, a substituição do reenvio integral do histórico pela combinação entre Memória Estável em \textit{Context Caching}, Memória Dinâmica e Janela Deslizante estabilizou o volume de \textit{tokens} transmitidos e permitiu o controle previsível de custos em sessões de longa duração.
    \item \textbf{Ampliação da Autonomia do Usuário (\textit{Human-in-the-Loop}) e dos Formatos de Entrega:} Enquanto na validação inicial de \citeonline{junior2025} os participantes apontaram variações na percepção de controle sobre o fluxo conversacional e na necessidade de ajustes manuais posteriores nas histórias exportadas em planilha (\texttt{.xlsx}), a nova interface incorporou mecanismos diretos de intervenção humana durante a própria entrevista. O estudante passou a dispor de edição manual de qualquer atributo dos oito tipos de artefatos em janela modal, exclusão lógica e reversão atômica do último turno de conversa (``Desfazer''), além da exportação consolidada do Documento de Requisitos de Software (DRS) completo em PDF, Markdown e JSON.
    \item \textbf{Governança Institucional e Ampliação do Cenário de Aplicação:} Por fim, a introdução do Painel de Administração dotou a ferramenta de controles operacionais necessários para o ambiente universitário, incluindo autenticação institucional Microsoft, restrição de uso às janelas de horário de aula, calibração dinâmica de versões do \textit{System Prompt} sem necessidade de reinicialização do servidor e auditoria de telemetria e \textit{logs}. Sob o ponto de vista operacional, essa infraestrutura possibilitou submeter a nova versão a testes automatizados de carga e a uma aplicação presencial em maior escala de participantes, passando de sessões individuais agendadas no estudo precursor ($N = 8$ respondentes) para duas turmas em laboratório (78 participantes no total, distribuídos em 7 equipes simultâneas por horário e $N = 65$ respondentes individuais), ainda que sob delineamento distinto e não diretamente comparável ao da avaliação anterior.
\end{enumerate}

\section{Discussão e Contribuições da Pesquisa}
\label{sec:contribuicoes}

Os resultados empíricos e operacionais consolidados neste estudo permitem analisar o comportamento da versão evoluída do ElicitBot frente aos fundamentos de Engenharia de Requisitos, o paradigma da Educação Inteligente (\textit{Smart Education}) e a arquitetura de agentes baseados em Modelos de Linguagem de Grande Porte (LLMs) apresentados no Capítulo \ref{cap:fundamentacao}. A seguir, é discutido como os achados da validação dialogam com a base teórica da área e a sistematização das contribuições científicas e práticas decorrentes da investigação.

\subsection{Confronto dos Resultados com a Literatura}

Na literatura de Engenharia de Requisitos, \citeonline{laplante2022} destacam que a atividade de elicitação não consiste em uma simples coleta de dados evidentes, mas em um processo investigativo rigoroso voltado a extrair necessidades, descobrir partes interessadas ocultas, refinar declarações iniciais vagas ou contraditórias e assegurar a rastreabilidade do escopo contra expansões descontroladas (\textit{scope creep}). De forma convergente, \citeonline{sommerville2016} ressalta que os \textit{stakeholders} frequentemente enfrentam dificuldades para articular o que desejam, tornando indispensável explorar o valor de negócio (o ``porquê'' de cada demanda) e validar se o sistema especificado corresponde à real necessidade do domínio. No contexto acadêmico, contudo, \citeonline{lojo2025} apontam que estudantes em formação apresentam dificuldades recorrentes para transpor esses fundamentos teóricos para a prática.

Os indicadores obtidos na validação em campo demonstram que a transição do ElicitBot para uma Máquina de Estados de 12 etapas atuou diretamente sobre esse desafio formativo. Ao superar o comportamento das versões preliminares (que iniciavam o diálogo solicitando diretamente funcionalidades), a nova arquitetura guiou os estudantes pelas etapas de contexto, problemas, partes interessadas, necessidades, objetivos e restrições antes da concepção técnica. Os índices de percepção reportados pelos estudantes, com 93,8\% de aprovação positiva em \textit{Investigação do Problema Real} ($\mu = 4,68$), 93,8\% em \textit{Compreensão da Dor Real} ($\mu = 4,52$), 100,0\% em \textit{Adesão ao Fluxo Lógico} ($\mu = 4,78$) e 96,6\% em \textit{Rastreabilidade de Requisitos} ($\mu = 4,60$), indicam que a condução implementada foi percebida em alinhamento com o ciclo investigativo, analítico e rastreável recomendado por \citeonline{sommerville2016} e \citeonline{laplante2022}.

Sob a ótica da Educação Inteligente (\textit{Smart Education}), os achados corroboram as proposições de \citeonline{zhu2016}, \citeonline{demir2021} e \citeonline{almogren2024} sobre a integração entre infraestrutura computacional adaptativa e metodologias ativas (\textit{Learning by Doing}). Conforme discutido por \citeonline{ataei2024}, ambientes de simulação funcionam como um espaço seguro para experimentação e erro, enquanto \citeonline{essel2022} e \citeonline{soliman2025} enfatizam o papel dos agentes conversacionais na superação das barreiras de escala impostas pela proporção numérica entre docentes e alunos. A aplicação prática do ElicitBot em duas turmas de horários distintos (78 estudantes presentes no total, distribuídos em 14 equipes com 7 sessões simultâneas por turma) ilustrou essa dinâmica em laboratório, posicionando o estudante como agente ativo e colaborativo na tomada de decisão e o professor como orientador metodológico e suporte tecnológico \cite{demir2021}.

No plano de orquestração de Modelos de Linguagem de Grande Porte (LLMs), a pesquisa dialoga com o alerta de \citeonline{ataei2024} e \citeonline{neumann2025} sobre a necessidade de mecanismos estritos de direcionamento para evitar respostas excessivamente genéricas. Especificamente quanto ao gerenciamento de janelas extensas de contexto, os dados operacionais dialogam com as conclusões de \citeonline{liu2023} sobre a degradação de atenção em históricos longos (\textit{Lost in the Middle}): mesmo em turnos que atingiram até 58.422 \textit{tokens} de entrada devido à colagem de textos prévios, a separação entre a Memória Estável (entidades confirmadas persistidas no banco) e a Janela Deslizante foi acompanhada por elevados índices de aprovação percebida em \textit{Adesão às Restrições} (96,8\%, $\mu = 4,63$) e \textit{Alinhamento da Solução} (95,3\%, $\mu = 4,53$). Adicionalmente, a incorporação da edição manual em modal e da reversão de turnos respondeu diretamente às observações de \citeonline{motger2021} e aos achados da validação inicial de \citeonline{junior2025}, na qual identificou-se variação na percepção de controle do usuário e necessidade de ajustes manuais posteriores nas histórias geradas.

\subsection{Contribuições Científicas e Acadêmicas}

Em alinhamento com o Ciclo de Rigor da \textit{Design Science Research} \cite{hevner2004, hevner2007}, esta investigação adiciona à base de conhecimento de Engenharia de Software e \textit{Smart Education} quatro contribuições principais:

\begin{enumerate}
    \item \textbf{Arquitetura Orientada a Estados e Validação Híbrida para LLMs em Engenharia de Requisitos:} A proposição de um modelo arquitetural que substitui o paradigma de conversação livre (\textit{Text-In, Text-Out}) por uma Máquina de Estados Finita de 12 etapas acoplada a um Motor de Operações transacional (\texttt{criar}, \texttt{atualizar}, \texttt{confirmar}, \texttt{validar} e \texttt{descartar}). Demonstrou-se como combinar a capacidade interpretativa de LLMs (via Saídas Estruturadas em JSON) com regras determinísticas de servidor para garantir integridade referencial e rastreabilidade entre problemas, necessidades, objetivos, restrições, soluções e requisitos funcionais.
    \item \textbf{Modelagem de \textit{Cloud Economics} e Gerenciamento Bipartido de Memória:} A formulação matemática e a validação empírica de uma estratégia de controle de contexto que divide os dados do projeto entre Memória Estável (congelada em \textit{Context Caching} na nuvem com desconto tarifário de 90\%) e Memória Dinâmica (artefatos em edição e janela deslizante). Esse modelo oferece um referencial replicável para pesquisas que buscam mitigar o crescimento quadrático de \textit{tokens} (``Efeito Avalanche'') em aplicações conversacionais de múltiplos turnos.
    \item \textbf{Evidências Empíricas em Escala de Turma com Triangulação de Telemetria:} A ampliação do corpo de evidências empíricas sobre o uso de LLMs no ensino de Engenharia de Requisitos \cite{khan2024, urquiza2025}, disponibilizando uma análise baseada em $N = 65$ avaliações discentes ($\alpha = 0,877$) cruzadas com 519 registros reais de telemetria de banco de dados. A identificação e documentação do comportamento de transposição acelerada de projetos pré-concebidos (\textit{Information Dumping}) fornece subsídios metodológicos relevantes para o desenho de futuros experimentos em Engenharia de Software Educacional.
    \item \textbf{Reprodutibilidade Científica:} A disponibilização pública dos artefatos de pesquisa, incluindo o histórico evolutivo dos \textit{System Prompts} (\texttt{v0.1} a \texttt{v1.0}), a memória de cálculo algébrica de custos, os \textit{scripts} de teste de carga em \textit{Locust} e as bases anonimizadas de respostas e telemetria, permitindo a auditoria independente e a replicação do estudo por outros pesquisadores.
\end{enumerate}

\subsection{Contribuições Práticas e Educacionais}

No âmbito do Ciclo de Relevância da DSR \cite{hevner2007}, voltado à solução concreta do problema vivenciado no ambiente acadêmico da Universidade de Brasília (UnB), destacam-se três contribuições práticas:

\begin{enumerate}
    \item \textbf{Entrega de Plataforma Educacional Persistente e Controlada:} A consolidação e implantação, em servidor institucional do grupo de pesquisa (CEDIS/UnB), de uma aplicação web completa, conteinerizada e resistente a quedas de conexão local. O registro de 0\% de erros nos testes de carga de até 50 usuários virtuais simultâneos e os 90,8\% de aprovação percebida pelos estudantes na validação presencial oferecem indícios favoráveis para a adoção controlada da ferramenta em semestres subsequentes da disciplina de Requisitos de Software.
    \item \textbf{Evidências de Viabilidade Financeira para Instituições Públicas de Ensino:} A demonstração empírica, nas condições das duas sessões avaliadas, de que o uso de LLMs comerciais em turmas universitárias pode ser operado a um custo faturado de R\$ 20,84 para o atendimento das duas turmas (R\$ 1,49 por equipe ou R\$ 0,27 por estudante presente), valor compatível com programas de créditos acadêmicos ou recursos de grupos de pesquisa.
    \item \textbf{Ferramental de Governança Docente e \textit{LLMOps}:} A entrega de um Painel de Administração integrado que permite aos professores gerenciar autonomamente a lista de alunos autorizados (\textit{whitelist}), restringir o consumo da API aos horários de aula, monitorar \textit{logs} de erro em tempo real, exportar relatórios de telemetria em CSV e calibrar novas versões do \textit{System Prompt} diretamente pela interface gráfica sem necessidade de intervenção na infraestrutura do servidor.
\end{enumerate}

\section{Limitações do Estudo e Ameaças à Validade}
\label{sec:limitacoes_estudo}

Embora os resultados empíricos e os testes de carga ofereçam indícios favoráveis quanto à viabilidade técnica, operacional e pedagógica da versão evoluída do ElicitBot nas condições avaliadas, a interpretação das conclusões deste trabalho deve considerar as limitações metodológicas, operacionais e de escopo inerentes ao delineamento adotado. A seguir, discutem-se as quatro principais frentes de limitação e ameaças à validade identificadas na pesquisa, explicitando o alcance dos procedimentos adotados para interpretar os dados sem pressupor a eliminação dos vieses subjacentes.

\subsection{Viés de Transposição de Escopo Pré-Concebido e Compressão Longitudinal}

A principal ameaça à validade ecológica do experimento decorreu do fato de as 14 equipes participantes utilizarem projetos semestrais cujo levantamento tradicional de requisitos já havia sido elaborado previamente na disciplina. Em vez de vivenciarem o ciclo de elicitação como um processo de descoberta gradual ao longo de várias semanas, os grupos utilizaram a ferramenta para transpor e validar especificações pré-existentes dentro de uma única sessão de 90 minutos (média de 37,1 turnos e 9,6 recriações de \textit{cache} por equipe).

Essa condição experimental induziu dois comportamentos atípicos: a aceleração das confirmações nas etapas iniciais de domínio (refletida na mediana de 19,0 \textit{tokens} nas mensagens dos estudantes) e o despejo prematuro de extensos blocos de documentação copiados para o \textit{chat} (\textit{information dumping}, com mensagens atingindo até 39.054 \textit{tokens}), o que gerou demandas por processamento em lote (\texttt{LA-06}) e por ingestão direta de arquivos (\texttt{TF-01}), além de acentuar a percepção de perda de detalhes periféricos após o giro da janela deslizante. A triangulação entre as respostas do questionário e os 519 registros de telemetria da entidade \texttt{MetricaLLM} auxiliou a contextualizar e interpretar o impacto dessa compressão sobre o consumo de \textit{tokens}, o armazenamento temporal de \textit{cache} (\textit{Storage TTL}) e a retenção de contexto, mas não elimina a diferença entre esse cenário e um uso longitudinal orgânico iniciado do zero.

\subsection{Dinâmica em Grupo, Vínculo com a Disciplina e Respostas Ausentes}

No âmbito da validade interna e de construção, diversos fatores associados ao protocolo de aplicação introduzem potenciais vieses nas respostas coletadas ($N = 65$). Primeiramente, a organização didática das turmas em grupos de 5 a 6 integrantes e a alocação de uma conta compartilhada por equipe geram um \textbf{efeito de avaliação em grupo}: como os integrantes de uma mesma equipe compartilharam a mesma sessão, a mesma tela projetada e os mesmos incidentes operacionais, suas avaliações individuais não são estatisticamente independentes entre si, além de o tempo de interação direta com o teclado não ter sido idêntico entre todos os membros, mesmo com o revezamento adotado.

Em segundo lugar, destaca-se o \textbf{efeito do vínculo com a disciplina}: por se tratar de uma atividade conduzida em turmas regulares da disciplina ministrada pelo orientador desta pesquisa, existe o risco de vieses de desejabilidade social, conformidade acadêmica ou entusiasmo inicial frente ao uso de Inteligência Artificial \cite{junior2025}. Em terceiro lugar, registra-se o impacto das \textbf{respostas ausentes (\textit{missing data})}: além de 13 dos 78 estudantes presentes não terem respondido ao formulário (taxa de retorno de 83,3\%), observou-se redução do número de respondentes nos itens das etapas finais ($N = 55$ a $61$ na Dimensão D4, e $n = 52$ respostas integrais no cálculo do Alfa de Cronbach), o que pode introduzir viés de autosseleção caso estudantes que enfrentaram maiores dificuldades tenham optado por se abster.

Embora a garantia de anonimato nas respostas assíncronas, o cálculo exploratório do Alfa de Cronbach ($\alpha = 0,877$), a triangulação com relatos abertos e telemetria e a Análise de Sensibilidade (variação máxima de $\pm 0,02$ ponto na Dimensão D4 ao filtrar participantes que alcançaram a Fase 10) contribuam para interpretar a consistência dos dados, tais procedimentos \textbf{não eliminam} os vieses de grupo, de vínculo institucional ou de não resposta, devendo os resultados ser lidos sob essas ressalvas.

\subsection{Restrições do Modelo Leve, Escopo Funcional e Versão Corrigida Não Reavaliada}

Sob a perspectiva técnica e de construção do artefato, a escolha do modelo \textit{Gemini 3.1 Flash-Lite} aliado à Janela Deslizante de 20 mensagens (10 turnos) representou um compromisso deliberado entre latência (mediana de 5.143 ms) e custo (R\$ 20,84 nas duas turmas). Como contrapartida, modelos compactos apresentam variações não determinísticas e menor capacidade de atenção seletiva diante de entradas longas com múltiplos contextos simultâneos (\texttt{LA-03}), o que contribuiu para a menor média relativa observada em \textit{Retenção de Contexto} ($\mu = 4,09$). Além disso, a telemetria persistida em \texttt{MetricaLLM} contabilizou exclusivamente os 519 turnos efetivados com sucesso, sem quantificar exaustivamente eventuais requisições revertidas ou com falha nos \textit{logs} rotacionados.

No plano funcional, o escopo permaneceu restrito à geração concisa de artefatos baseada apenas nos dados informados pelo usuário (\texttt{LA-02}) e ao padrão Connextra de Histórias de Usuário (\texttt{LA-05}), postergando para trabalhos futuros outros \textit{frameworks} como o FDD (\texttt{TF-04}), Requisitos Não Funcionais (\textbf{US2.06}), entrada por voz (\textbf{US2.07} e \texttt{TF-02}), exportação para plataformas externas (\textbf{US3.03}) e a reversão de edições parciais ou de mudanças de fase no comando ``Desfazer''. Por fim, constitui limitação relevante o fato de a \textbf{versão corrigida após a validação em campo não ter sido reavaliada com estudantes}: as intervenções imediatas \texttt{CI-01} a \texttt{CI-04} (incluindo a correção da falha de concorrência na Fase 11 e a reformulação do campo de entrada) foram verificadas apenas por testes funcionais manuais do autor após a reimplantação, de modo que os resultados do questionário referem-se exclusivamente à versão utilizada na sessão de 22 de outubro de 2026 e não comprovam empiricamente o desempenho da versão refinada em sala de aula.

\subsection{Recorte Amostral, Institucional e Delimitação das Evidências de Aprendizagem}

Quanto à validade externa e de conclusão, a aplicação em campo restringiu-se a estudantes de graduação em Engenharia de Software de uma única instituição pública (UnB, Campus Gama / FCTE) em uma única janela letiva, sem grupo de controle submetido paralelamente a métodos tradicionais e sem acompanhamento longitudinal ao longo do semestre.

Além disso, delimita-se estritamente o que as evidências coletadas permitem afirmar sobre aprendizagem e qualidade dos entregáveis. Este estudo \textbf{não realizou uma avaliação independente dos artefatos produzidos} pelas 14 equipes por meio de rubricas docentes ou revisores externos, tampouco aplicou pré-testes e pós-testes formais para mensurar ganho cognitivo. As evidências disponíveis restringem-se à percepção autodeclarada pelos estudantes no questionário ($N = 65$), aos relatos qualitativos abertos, aos registros de telemetria operacional e às validações estruturais internas executadas automaticamente pelo próprio sistema (contratos de esquema Pydantic e regras determinísticas de integridade referencial no \textit{backend}). Consequentemente, os dados suportam inferências sobre aceitação, utilidade percebida, aderência operacional ao fluxo e viabilidade técnica nas condições estudadas, mas não permitem comprovar diretamente a qualidade semântica intrínseca dos requisitos gerados nem ganhos objetivos de aprendizagem. A disponibilização aberta do código-fonte, dos \textit{prompts} e das bases anonimizadas visa viabilizar investigações futuras com avaliação independente de artefatos e delineamentos quase-experimentais.

Consolidada a análise crítica dos resultados, a triangulação metodológica, o comparativo evolutivo entre versões, as contribuições e as limitações da pesquisa, o capítulo final sintetiza as conclusões gerais do trabalho e projeta os desdobramentos para investigações futuras.